\chapter{Компьютер из \nt{СЕРО}-нейронов}
\label{sec:serocomputer}

\section{Чем \nt{СЕРО}-нейрон отличается от \nt{ГЛУТ}}

В главе~\ref{sec:glutcomputer} мы собрали компьютер из \nt{ГЛУТ}-нейронов и упёрлись в одно ограничение: все веса положительны, поэтому импульс нельзя погасить импульсом. Давайте повторим опыт с другим типом. Возьмём сколько угодно \nt{СЕРО}-нейронов и только их. Что изменится?

С точки зрения инженера у \nt{СЕРО}-нейрона четыре важных отличия.

\emph{Первое: знак выбирает получатель.} В главе~\ref{sec:neurontypes} мы видели, что знак действия медиатора задаёт рецептор (утверждение~\ref{prop:receptor}). У серотонина есть рецепторы обоих знаков: одни облегчают возбуждение получателя, другие затрудняют. Значит, в сети из \nt{СЕРО}-нейронов веса бывают и положительными, и отрицательными. Правило~\eqref{eq:dalesign} здесь поворачивается на девяносто градусов: знак определяется не отправителем, а тем, какой рецептор поставил получатель.

\emph{Второе: \nt{СЕРО}-нейрон работает сам.} Серотониновые нейроны выдают импульсы ровно, несколько раз в секунду, даже когда на них ничего не действует. Это ритмоводитель: у него есть внутренний источник возбуждения, который в модели выглядит как постоянное слагаемое $b_i$ в сумме.

\emph{Третье: он медленный.} Почти все рецепторы серотонина метаботропные: они не открывают канал сами, а запускают внутри клетки цепочку реакций. Ответ приходит через сотни миллисекунд и длится секунды. Шаг времени в нашей модели будет теперь порядка секунды, а не миллисекунды.

\emph{Четвёртое: он выбрасывает медиатор не в провод, а в объём.} Серотонин часто выходит не в узкую щель, а в межклеточное пространство и расплывается вокруг. Получатель чувствует концентрацию и не может знать, от какого отправителя пришла молекула.

Четвёртое отличие мы пока отложим до раздела~\ref{sec:volume} и будем считать, что выход каждого нейрона доходит только до тех, кому он предназначен. Первые три меняют правило~\eqref{eq:glutneuron} так:
\begin{empheq}[box=\keybox]{equation}
y_i(t+1) \;=\; \Theta\Bigl(\sum_j w_{ij}\,y_j(t) \;+\; b_i \;-\; \theta_i\Bigr), \qquad w_{ij}\gtrless 0,\quad b_i\ge 0 .
\label{eq:seroneuron}
\end{empheq}
Веса теперь любого знака, а нейрон с $b_i\ge\theta_i$ срабатывает каждый шаг, пока его не затормозят. Один шаг --- около секунды.

На рисунках тормозящую связь будем рисовать линией с поперечной чертой на конце, а ритмоводитель --- двойным кружком.

\section{НЕ за один нейрон}

Сделаем НЕ. Возьмём ритмоводитель с $b=\theta=1$ и подадим на него вход $x$ с весом $-1$. Если вход молчит, сумма равна $b=1$, нейрон срабатывает. Если вход срабатывает, сумма равна нулю, нейрон молчит. Это НЕ, и на него ушёл один нейрон (рис.~\ref{fig:serogates}). Утверждение~\ref{prop:monotone} здесь неприменимо: оно требовало $w_{ij}\ge0$.

Подадим на тот же ритмоводитель два тормозящих входа. Он молчит, если сработал хотя бы один из них, и срабатывает, только если оба молчат. Это ИЛИ-НЕ. А ИЛИ-НЕ --- элемент особый: из него одного можно собрать любую логическую схему. НЕ --- это ИЛИ-НЕ с одним входом, ИЛИ --- ИЛИ-НЕ, за которым стоит НЕ, а И получается по правилу де Моргана: $x\wedge y=\overline{\bar x\vee\bar y}$.

\begin{figure}[t]
\centering
\begin{tikzpicture}[>=Stealth,
  nrn/.style={circle,draw,very thick,fill=blue!6,minimum size=0.95cm,inner sep=0pt},
  pace/.style={nrn,double,double distance=1.2pt},
  inh/.style={-{Bar[width=7pt]},thick,red!70!black}]
\node[pace] (N1) at (0,0) {$1$};
\draw[inh] (-1.6,0) node[left,black] {$x$} -- (N1);
\draw[->,very thick,red!70!black] (N1) -- (1.3,0) node[right,black] {$\bar x$};
\node[below=0.55cm] at (0,0) {\small НЕ};
\node[pace] (N2) at (5,0) {$1$};
\draw[inh] (3.4,0.45) node[left,black] {$x$} -- (N2);
\draw[inh] (3.4,-0.45) node[left,black] {$y$} -- (N2);
\draw[->,very thick,red!70!black] (N2) -- (6.3,0) node[right,black] {$\overline{x\vee y}$};
\node[below=0.55cm] at (5,0) {\small ИЛИ-НЕ};
\end{tikzpicture}
\caption{Логика из \nt{СЕРО}-нейронов. Двойной кружок --- ритмоводитель: $b=\theta=1$, он срабатывает каждый шаг, пока его не затормозят. Линия с чертой --- тормозящая связь с весом $-1$. Слева --- НЕ, справа --- ИЛИ-НЕ; и то и другое --- один нейрон.}
\label{fig:serogates}
\end{figure}

\begin{keyblock}
\begin{proposition}[Полнота \nt{СЕРО}-логики]
\label{prop:serocomplete}
Сеть из нейронов~\eqref{eq:seroneuron}, в которой есть ритмоводители и тормозящие связи, может вычислить любую логическую функцию. Парафазный код для этого не нужен: каждый бит передаётся одним проводом.
\end{proposition}
\end{keyblock}

Сравните с главой~\ref{sec:glutcomputer}. Там ради НЕ пришлось удвоить все провода и почти все нейроны. Здесь НЕ стоит один нейрон, а константа «единица» даётся даром: это любой ритмоводитель без тормозящих входов.

\section{Сумматор из двух нейронов}

Отрицательные веса делают пороговый нейрон ещё мощнее. Вспомним полный сумматор: входы $a,b,c$, число единиц $n=a+b+c$, перенос $c'=1$ при $n\ge2$, сумма $s=1$ при нечётном $n$. Перенос --- по-прежнему один нейрон большинства. А сумму теперь тоже можно получить одним нейроном, если дать ему вход от нейрона переноса с весом $-2$:
\begin{empheq}[box=\keybox]{equation}
s \;=\; \Theta\bigl(\,a+b+c \;-\; 2\,c' \;-\; 1\,\bigr), \qquad c' = \Theta\bigl(\,a+b+c-2\,\bigr).
\label{eq:serosum}
\end{empheq}
Проверим. При $n=0$: $0-0=0<1$, молчит. При $n=1$: $1-0=1$, срабатывает. При $n=2$: $2-2=0$, молчит. При $n=3$: $3-2=1$, срабатывает. Всё верно (рис.~\ref{fig:seroadder}).

\begin{figure}[t]
\centering
\begin{tikzpicture}[>=Stealth,
  nrn/.style={circle,draw,very thick,fill=blue!6,minimum size=0.95cm,inner sep=0pt},
  inh/.style={-{Bar[width=7pt]},thick,red!70!black}]
\node[nrn] (C) at (0,-1.1) {$2$};
\node[nrn,fill=red!8] (S) at (3,0.6) {$1$};
\node[font=\small] (in) at (-2.6,0.6) {$a,b,c$};
\draw[->,thick,blue!60!black] (in.east) -- (S);
\draw[->,thick,blue!60!black] (in.south east) -- (C);
\draw[inh] (C) -- node[below right=-2pt] {\scriptsize $-2$} (S);
\draw[->,very thick,red!70!black] (S) -- (4.6,0.6) node[right,black] {$s$};
\draw[->,very thick,red!70!black] (C) -- (4.6,-1.1) node[right,black] {$c'$};
\end{tikzpicture}
\caption{Полный сумматор из двух \nt{СЕРО}-нейронов. Нижний нейрон --- большинство, он даёт перенос $c'$. Верхний складывает входы и получает от переноса торможение с весом $-2$; его порог $1$. Сравните с восемью нейронами на рис.~\ref{fig:fulladder}.}
\label{fig:seroadder}
\end{figure}

Два нейрона на разряд вместо восьми. Цепочка переносов проходит, как и раньше, один разряд за шаг. Мы проверили четырёхразрядный сумматор на модели для всех $256$ пар слагаемых. При сложении $7+1$ выход последовательно показывает $7,6,4$ и на четвёртом шаге устанавливается на правильном ответе $8$. Промежуточные неверные значения --- те же гонки переноса, что и в главе~\ref{sec:glutcomputer}.

\section{Память со сбросом}

Ячейка памяти тоже становится проще. Раньше петле хранения нужен был отдельный сигнал $h$, потому что погасить петлю можно было, только убрав его. Теперь петлю можно погасить торможением. Нейрон $H$ с весами $+1$ от $q$ и $-1$ от сигнала записи $e$ и порогом $1$ вычисляет «$q$ и не $e$»: он поддерживает петлю, пока запись не разрешена. Нейрон $W$ с порогом $2$ вычисляет «$d$ и $e$», а нейрон $q$ --- ИЛИ от $H$ и $W$. Три нейрона на бит вместо шести, и сигнал хранения больше не нужен: когда $e$ открывает запись, он же и гасит старое значение.

Можно ли обойтись одним нейроном? Нет. Рассуждение из главы~\ref{sec:glutcomputer}, где мы сложили четыре неравенства и пришли к противоречию, нигде не использовало знак весов. Оно показывает, что переключатель «хранить или записать» не может сделать один пороговый нейрон ни с какими весами. Это свойство самого порогового элемента, а не медиатора.

\section{Такт, который запускается сам}
\label{sec:johnson}

Самое интересное отличие --- в тактовом генераторе. В главе~\ref{sec:glutcomputer} кольцо из \nt{ГЛУТ}-нейронов надо было запустить внешним импульсом, а лишний импульс, попавший в кольцо, оставался там навсегда. Теперь у нас есть НЕ.

Соединим $K$ нейронов в цепочку, как раньше, но на вход первого подадим \emph{инвертированный} выход последнего. Такое кольцо в электронике называют счётчиком Джонсона. Начнём с тишины: все $K$ нейронов молчат. Первый нейрон видит молчание последнего, инвертирует его и срабатывает. Дальше единицы по одной заполняют цепочку, пока не сработают все $K$. Тогда первый нейрон видит единицу, замолкает, и цепочку по одному заполняют нули. Через $2K$ шагов всё повторяется. При $K=8$:
\[
\texttt{10000000}\to\texttt{11000000}\to\dots\to\texttt{11111111}\to\texttt{01111111}\to\dots\to\texttt{00000000}\to\texttt{10000000}.
\]
Из восьми нейронов получилось шестнадцать фаз, и генератор \emph{запустился сам}, из полной тишины.

А что с лишним импульсом? Простое кольцо Джонсона от него не лечится: если перевернуть один бит, неправильный узор так и будет бегать по кругу. Мы проверили все $256$ состояний восьми нейронов: из $240$ неправильных ни одно не возвращается к правильному циклу. Но теперь мы можем это исправить, потому что умеем гасить импульс импульсом. Дадим первому нейрону кольца, $J_0$, вход от самого себя с весом $+1$, тормозящие входы от четвёртого и восьмого нейронов с весом $-2$ и постоянное слагаемое $b=2$ при пороге $1$ (рис.~\ref{fig:johnson}):
\begin{equation}
J_0(t+1) \;=\; \Theta\bigl(\,J_0 - 2J_3 - 2J_7 + 2 - 1\,\bigr).
\label{eq:johnsonfix}
\end{equation}
Правильный цикл от этого не меняется, а из \emph{любого} из $256$ состояний кольцо не более чем за $13$ шагов возвращается к нему. Такой генератор сам запускается и сам исправляет сбои. Для \nt{ГЛУТ}-кольца это было невозможно в принципе: утверждение~\ref{prop:monotone} запрещало гасить лишний импульс.

\begin{figure}[t]
\centering
\begin{tikzpicture}[>=Stealth,
  nrn/.style={circle,draw,very thick,fill=blue!6,minimum size=0.8cm,inner sep=0pt},
  pace/.style={nrn,double,double distance=1.2pt},
  inh/.style={-{Bar[width=7pt]},thick,red!70!black}]
\node[pace] (J0) at (0,0) {\small $J_0$};
\foreach \k in {1,...,7}{\node[nrn] (J\k) at (1.45*\k,0) {\small $J_\k$};}
\foreach \k [evaluate=\k as \p using int(\k-1)] in {1,...,7}{\draw[->,thick] (J\p) -- (J\k);}
\draw[inh] (J7.south) -- ++(0,-0.9) -| node[pos=0.25,below] {\scriptsize $-2$} (J0.south);
\draw[inh] (J3.north) -- ++(0,0.8) -| node[pos=0.3,above] {\scriptsize $-2$} (J0.north east);
\draw[->,thick] (J0.north west) .. controls (-0.9,1.1) and (-1.1,-0.2) .. node[left] {\scriptsize $+1$} (J0.west);
\end{tikzpicture}
\caption{Самоисправляющийся тактовый генератор из восьми \nt{СЕРО}-нейронов. Каждый нейрон повторяет предыдущий. Первый, $J_0$, --- ритмоводитель с $b=2$ и порогом $1$; его тормозят четвёртый и последний нейроны, а сам он поддерживает себя. Кольцо выдаёт $16$ фаз, запускается из тишины и из любого сбоя возвращается к правильному циклу не более чем за $13$ шагов.}
\label{fig:johnson}
\end{figure}

Сигналы записи для регистров теперь тоже дешевле. Нужная пара фаз кольца выделяется одним нейроном с одним возбуждающим и одним тормозящим входом: «нейрон $J_i$ сработал, а нейрон $J_j$ ещё нет». Отдельные сигналы хранения не нужны вовсе.

\section{Тот же счётчик}

Соберём тот же четырёхразрядный счётчик, что и в главе~\ref{sec:glutcomputer}: регистр $A$, сумматор, прибавляющий единицу, регистр $T$ и тактовый генератор. Архитектура та же (рис.~\ref{fig:counter}), меняются только детали.

\begin{table}[t]
\centering
\small
\begin{tabular}{lrr}
\toprule
Узел & \nt{ГЛУТ} & \nt{СЕРО} \\
\midrule
тактовый генератор & 16 & 8 \\
сигналы управления & 4 & 2 \\
источник единицы & 1 & 1 \\
сумматор, 4 разряда & 32 & 8 \\
регистры $A$ и $T$ & 48 & 24 \\
\midrule
всего нейронов & 101 & 43 \\
\bottomrule
\end{tabular}
\caption{Четырёхразрядный счётчик из \nt{ГЛУТ}-нейронов (глава~\ref{sec:glutcomputer}) и из \nt{СЕРО}-нейронов.}
\label{tab:serocounter}
\end{table}

Получилось $43$ нейрона вместо $101$ (таблица~\ref{tab:serocounter}). Мы промоделировали эту сеть по правилу~\eqref{eq:seroneuron}, начав с полной тишины: ни одного импульса ни в кольце, ни в регистрах. Ритмоводители запустились сами, и через один такт регистр $A$ начал показывать $1,2,3,\dots,15,0,1,\dots$ --- по одному числу за $16$ шагов.

Но шаг теперь около секунды. Значит, счётчик прибавляет единицу раз в шестнадцать секунд, а сумматор на $32$ разряда устанавливается за полминуты. Сведём сравнение в одну таблицу.

\begin{table}[t]
\centering
\small
\begin{tabular}{>{\raggedright}p{4.2cm}>{\raggedright}p{4.0cm}>{\raggedright\arraybackslash}p{4.0cm}}
\toprule
& \nt{ГЛУТ} & \nt{СЕРО} \\
\midrule
знак весов & только $+$ & $+$ и $-$, выбирает получатель \\
НЕ & перестановка проводов, все провода удвоены & один нейрон \\
полный сумматор & 8 нейронов & 2 нейрона \\
бит памяти & 6 нейронов & 3 нейрона \\
запуск & нужен внешний импульс & из тишины \\
лишний импульс в такте & остаётся навсегда & исправляется за 13 шагов \\
шаг времени & $\sim1$~мс & $\sim1$~с \\
один отсчёт счётчика & 16~мс & 16~с \\
\bottomrule
\end{tabular}
\caption{Компьютер из \nt{ГЛУТ}-нейронов и компьютер из \nt{СЕРО}-нейронов.}
\label{tab:glutvssero}
\end{table}

По числу нейронов и по надёжности \nt{СЕРО}-компьютер выигрывает с большим запасом. По скорости он проигрывает в тысячу раз. Но это ещё не главная его беда.

\section{Провод --- это объём}
\label{sec:volume}

Всё это время мы жульничали. Мы считали, что выход каждого \nt{СЕРО}-нейрона доходит только до тех получателей, которым он предназначен. Но серотонин выходит в межклеточное пространство и расплывается. Получатель чувствует концентрацию серотонина вокруг себя и не может отличить молекулу от одного отправителя от молекулы от другого. Всё, что выброшено поблизости, складывается.

Что это значит для нашей схемы? \emph{Провод --- это не волокно, а объём ткани.} Два сигнала остаются разными, только если их выбрасывают в области, разнесённые дальше, чем успевает расплыться медиатор. Это расстояние легко оценить. За время $\tau$ молекула с коэффициентом диффузии $D$ уходит на
\begin{empheq}[box=\keybox]{equation}
\ell \;\sim\; \sqrt{D\,\tau}\, .
\label{eq:diffusionlength}
\end{empheq}
Для небольшой молекулы в межклеточном пространстве мозга $D$ порядка $10^{-6}$~см$^2$/с; если сигнал живёт $\tau\sim1$~с, то $\ell\sim\sqrt{10^{-6}}$~см $=10$~мкм. Грубо, область одного провода --- шар в десяток-другой микрометров. Адрес в таком компьютере --- это \emph{место}: получатель узнаёт, от кого сигнал, только по тому, где он сам находится.

Нашему счётчику нужно $43$ таких отдельных области. Это выполнимо, но настоящие \nt{СЕРО}-нейроны устроены ровно наоборот. Вспомните таблицу~\ref{tab:types}: у каждого из них около $10^5$ мест выброса, раскиданных по огромным участкам мозга. Один нейрон --- это один «провод», протянутый через огромный объём, и «провода» разных \nt{СЕРО}-нейронов перекрываются и смешиваются. К тому же многие серотониновые нейроны работают почти в унисон. В итоге вся серотониновая система мозга --- это не сотни тысяч независимых проводов, а немного, может быть, всего несколько, медленных общих переменных.

\begin{keyblock}
\begin{proposition}[Число независимых проводов]
\label{prop:volumewires}
В сети с объёмной передачей число независимых сигналов ограничено не числом нейронов, а числом непересекающихся областей размером $\ell\sim\sqrt{D\tau}$, в которые они выбрасывают медиатор. Если выход каждого нейрона раскидан по большому объёму, вся сеть передаёт лишь несколько общих переменных.
\end{proposition}
\end{keyblock}

Вот и итог двух глав. Как логический элемент \nt{СЕРО}-нейрон богаче \nt{ГЛУТ}: у него есть оба знака и встроенный ритм, из него получается и НЕ, и самоисправляющийся такт. Но как система связи он --- широковещательная рассылка: медленная и почти без адресов. Настоящий \nt{СЕРО} мозга не пытается быть процессором. Он делает ровно то, для чего годится такая связь: медленно сообщает всему мозгу несколько общих чисел, о которых шла речь в главе~\ref{sec:neurontypes}. А быстрые адресные вычисления мозг поручает \nt{ГЛУТ}, дополняя его тормозящими \nt{ГАМК}-нейронами, которые умеют гасить импульс импульсом так же, как это делали наши тормозящие связи, но быстро и адресно.
